\begin{titlepage}
    % Strona tytułowa
    \vbox to\textheight{\hyphenpenalty=10000
    \begin{center}
        \begin{tabular}{p{107mm} p{9cm}}
            \begin{minipage}{9cm}
              \begin{center}
              Politechnika Warszawska \\
              Wydział Elektroniki i~Technik Informacyjnych \\
              Instytut Informatyki
              \end{center}
            \end{minipage}
            &
            \begin{minipage}{8cm}
            \begin{flushleft}
             \footnotesize
              Rok akademicki 2011/2012
            \vspace*{2.75\baselineskip}
            \end{flushleft}
            \end{minipage} \\
        \end{tabular}
        \vspace*{3.75\baselineskip}
        \par\vspace{\smallskipamount}
        \includegraphics[width=4cm]{build/pw.eps}\par
        \vspace*{2\baselineskip}{\LARGE Praca dyplomowa inżynierska\par}
        \vspace{3\baselineskip}{\LARGE\strut \theauthor\par}
        \vspace*{2\baselineskip}{\huge\bfseries \thetitle\par}

        \vspace*{7\baselineskip}
        \hfill\mbox{}\par\vspace*{\baselineskip}\noindent
        \begin{tabular}[b]{@{}p{3cm}@{\ }l@{}}
            {\large\hfill } & {\large }
        \end{tabular}
        \hfill
        \begin{tabular}[b]{@{}l@{}}
        Opiekun pracy: \\[\smallskipamount]
        {\large dr inż.\ Wojciech Zabołotny}
        \end{tabular}\par
        \vspace*{4\baselineskip}
    \begin{tabular}{p{\textwidth}}
    \begin{flushleft}
        \begin{minipage}{7cm}
        Ocena \dotfill
        \par\vspace{1.6\baselineskip}
        \dotfill
        \par\noindent
        \centerline{\footnotesize Podpis Przewodniczącego} \par
        \centerline{\footnotesize Komisji Egzaminu Dyplomowego}\par
        \end{minipage}
    \end{flushleft}
    \end{tabular}
    \end{center}}

    % Życiorys
    \newpage\thispagestyle{empty}
    \begin{tabular}{p{5cm} p{12cm}}
    \begin{minipage}{5cm}
    \center
    \IfFileExists{bsc-photo.eps}{
      \includegraphics[width=4.5cm]{bsc-photo.eps}
    }{
      \includegraphics[width=4.5cm]{build/photo.eps}
    }
    \end{minipage}
    &
    \begin{minipage}{12cm}
    \begin{flushleft}
    \par\noindent\vspace{1\baselineskip}
    \begin{tabular}[h]{l l}
    {\it Kierunek:}       & Informatyka \\[12pt]
    {\it Specjalność:}    & Inżynieria Systemów \\
                          & Informatycznych \\[12pt]
    {\it Data urodzenia:} & 31 sierpnia 1986~r. \\[12pt]
    {\it Data rozpoczęcia studiów:} & 1 lutego 2006 r. \\
    \end{tabular}
    \par\noindent\vspace{1\baselineskip}
    \end{flushleft}
    \end{minipage}
    \end{tabular}
    \vspace*{1\baselineskip}
    \begin{center}
        {\large\bfseries Życiorys}\par\bigskip
    \end{center}

    \IfFileExists{bsc-bio.tex}{\indent\input{bsc-bio}}{
      \begin{center}
        $\ddot\smile$
      \end{center}
    }
    \par
    \vspace{2\baselineskip}
    \hfill\parbox{15em}{{\small\dotfill}\\[-.3ex]
    \centerline{\footnotesize podpis studenta}}\par
    \vspace{3\baselineskip}
    \begin{center}
        {\large\bfseries Egzamin dyplomowy} \par\bigskip\bigskip
    \end{center}
    \par\noindent\vspace{1.5\baselineskip}
    Złożył egzamin dyplomowy w dn. \dotfill
    \par\noindent\vspace{1.5\baselineskip}
    Z wynikiem \dotfill
    \par\noindent\vspace{1.5\baselineskip}
    Ogólny wynik studiów \dotfill
    \par\noindent\vspace{1.5\baselineskip}
    Dodatkowe wnioski i uwagi Komisji \dotfill
    \par\noindent\vspace{1.5\baselineskip}
    \dotfill

    % Streszczenie
    \newpage\thispagestyle{empty}
    \vspace*{2\baselineskip}
    \begin{center}
        {\large\bfseries Streszczenie}\par\bigskip
    \end{center}

    {\itshape Wiele podzespołów komputera, a~szczególnie
      tzw.\ systemów wbudowanych, jest zazwyczaj podłączonych
      bezpośrednio do magistrali systemowej, przez co muszą operować
      adresami fizycznymi.  Równocześnie mechanizmy bezpośredniego
      dostępu do pamięci (\ang{Direct Memory Access}, \acc{DMA}) są
      ograniczone do transferów sekwencyjnych.  Stwarza to potrzebę,
      alokowania dużych ciągłych fizycznie obszarów pamięci do
      wykorzystania w~takich układach.

      Dotychczas stosowane w~systemach bazujących na jądrze Linux
      rozwiązania wiążą się z~rezerwowaniem dużego obszaru pamięci,
      który wyjęty spod kontroli Linuksa przestaje być użyteczny dla
      jądra i~w~rezultacie wykorzystywany nieefektywnie.

      Praca opisuje stworzony przeze mnie alokator pamięci ciągłej,
      \ang{Contiguous Memory Allocator} (\acc{CMA}), który rozwiązuje ten
      problem poprzez zastosowanie mechanizmu migracji, który pozwala
      przenosić zajęte strony i~w~ten sposób tworzyć długie sekwencje
      wolnych stron.}

    \vspace*{1\baselineskip}

    \noindent{\bf Słowa kluczowe}: {\itshape alokacja pamięci, Linux,
      systemy wbudowane.}
    \par
    \vspace{4\baselineskip}
    \begin{center}
        {\large\bfseries Abstract}\par\bigskip
    \end{center}
    \noindent{\bf Title}: {\itshape \theengishtitle}\par
    \vspace*{1\baselineskip}

    {\itshape Many computer components, especially in a~so called
      embedde system, are attached directly to the system bus and thus
      need to operate on physical addresses.  At the same time, Direct
      Memory Access (\acc{DMA}) is limited to sequential transfers only.
      This creates a need to allocate big physically contiguous memory
      buffers to be used with such components.

      Solutions previously used in Linux-based systems boil down to
      reserving big memory area which removed from Linux control
      cannot be used by the kernel and is used inefficiently.

      This work describes Contiguous Memory Allocator (\acc{CMA}) -- an
      allocator written by me which solves the problem by using
      migration, which makes it possible to move allocated pages and
      thus create a~long sequence of free pages.}
    \vspace*{1\baselineskip}

    \noindent{\bf Key words}: {\itshape memory allocation, Linux,
      embedded systems.}

\end{titlepage}
